\textbf{结论名称：} 标准差对线性变换的响应性质（平移不变性与缩放比例性）

\textbf{适用条件：} 对任意一组实数值数据 \(x_1, x_2, \dots, x_n\)（\(n \ge 2\)），令 \(y_i = a + kx_i\)（其中 \(a, k\) 为常数，\(k \neq 0\)）。结论也适用于总体标准差与样本标准差。

\textbf{变量说明：} 设原数据为 \(x_i\)，其标准差记为 \(s_x\)；变换后数据 \(y_i = a + k x_i\)，标准差记为 \(s_y\)。常数 \(a\) 为平移量，\(k\) 为缩放因子（可正可负）。

\textbf{核心公式：}
\[
\boxed{s_y = |k| \cdot s_x}
\]
特别地，当 \(k=0\) 时所有 \(y_i\) 相等，\(s_y = 0\)（与 \(0\cdot s_x=0\) 一致）。平移 \(a\) 不改变标准差：
\[
s_{x + a} = s_x
\]

\textbf{使用场景：} 数据单位换算（如摄氏度转华氏度）、数据标准化（z-score）、处理调整后的数据集时快速计算标准差；在回归分析或方差分析中理解尺度变换对离散程度的影响。

\textbf{简单例子：} 已知数据 \(x\)：\(2, 4, 6\)，其均值 \(\bar{x}=4\)，标准差
\[
s_x=\sqrt{\frac{(2-4)^2+(4-4)^2+(6-4)^2}{3}}=\sqrt{\frac{4+0+4}{3}}=\sqrt{\frac{8}{3}}\approx 1.633.
\]
令 \(y_i = 3 + 2x_i\)，得 \(y\)：\(7, 11, 15\)。均值 \(\bar{y}=3+2\cdot4=11\)，直接计算 \(s_y=\sqrt{\frac{(7-11)^2+(11-11)^2+(15-11)^2}{3}}=\sqrt{\frac{16+0+16}{3}}=\sqrt{\frac{32}{3}}\approx 3.266\)，恰好等于 \(|2|\cdot s_x = 2\sqrt{8/3} = \sqrt{32/3}\)，平移量 \(3\) 不影响标准差。